% Compile with XeLaTeX and a TeX distribution that includes ctex and Fandol.
% For an English-only resume, use \documentclass[11pt,a4paper]{article} and translate the section headings.
\documentclass[11pt,a4paper,scheme=plain,fontset=fandol]{ctexart}
\usepackage[margin=19mm]{geometry}
\usepackage{enumitem}
\usepackage[hidelinks]{hyperref}
\pagestyle{empty}
\setlength{\parindent}{0pt}
\setlength{\parskip}{3pt}
\setlist[itemize]{leftmargin=*, topsep=3pt, itemsep=2pt, parsep=0pt}

\begin{document}

{\LARGE\bfseries 候选人姓名}\par
\href{mailto:hello@example.com}{hello@example.com}
\enspace\textbar\enspace \href{https://example.com}{作品集}\par
学校 · 专业 · 毕业年份\par
开源贡献或代表奖项（按需保留）

\section*{工作经历}
\textbf{公司 · 岗位} \hfill YYYY.MM--至今\par
\textbf{项目名称 · 项目角色}
\begin{itemize}
  \item \textbf{情境}：服务对象遇到了什么问题。
  \item \textbf{任务}：项目目标与本人负责的交付。
  \item \textbf{行动}：本人做了什么，用了什么方法。
  \item \textbf{结果}：已验证的交付、采纳或指标。
\end{itemize}

\section*{实习经历}
\textbf{公司 · 岗位} \hfill YYYY.MM--YYYY.MM\par
\textbf{项目名称}
\begin{itemize}
  \item \textbf{情境}：服务对象遇到了什么问题。
  \item \textbf{任务}：项目目标与本人负责的交付。
  \item \textbf{行动}：本人做了什么，用了什么方法。
  \item \textbf{结果}：已验证的交付、采纳或指标。
\end{itemize}

\section*{AI 创业／独立项目}
\textbf{项目名称 · 项目角色} \hfill YYYY.MM--YYYY.MM
\begin{itemize}
  \item \textbf{情境}：服务对象遇到了什么问题。
  \item \textbf{任务}：项目目标与本人负责的交付。
  \item \textbf{行动}：本人做了什么，用了什么方法。
  \item \textbf{结果}：已验证的交付、采纳或指标。
\end{itemize}

\end{document}
